\section{True vertex-weight components and the whole formal Hilbert--Samuel bound}

\begin{definition}[The actual coefficientwise weighted component]
\label{hir:actual-weighted-component-definition}
\lean{Hironaka.HironakaActualWeightedVertexHS.actualWeightedComponent,Hironaka.HironakaActualWeightedVertexHS.actualWeightedComponent_coefficient}
\leanok
For any nonnegative coordinate weights $w$, degree $d$ and original
formal series $F$, retain exactly the coefficients with $w$-weight $d$
and set all other coefficients to zero. This defines a genuine formal
series, without positivity of the weights or finite-support assumptions.
\end{definition}
\begin{proof}
An original multivariate formal series is precisely a function from
actual finite exponents to complex coefficients. The displayed
coefficient rule is therefore an actual original series, and evaluation
at each exponent gives exactly the asserted coefficient formula. Zero
weights may retain infinitely many tangential coefficients; no analytic
convergence or finite-polynomial claim is made.
\end{proof}

\begin{theorem}[The component at a genuine initial exponent's weight retains that exact exponent]
\label{hir:actual-weighted-component-retains-initial-exponent}
\uses{hir:actual-weighted-component-definition}
\lean{Hironaka.HironakaActualWeightedVertexHS.actualWeightedComponent_ne_zero,Hironaka.HironakaActualWeightedVertexHS.actualWeightedComponent_initialExponent}
\leanok
Let $F$ be an actual nonzero formal series and $a$ its original least
nonzero exponent for the given genuine monomial order. The actual
component of $F$ at degree $w(a)$ is nonzero and has precisely that same
original initial exponent $a$.
\end{theorem}
\begin{proof}
Its coefficient at $a$ is the unchanged nonzero coefficient of $F$,
since $w(a)=w(a)$. Thus the component is nonzero. At every smaller
exponent the coefficient of $F$ was zero. The component either retains
that zero coefficient or replaces it by zero. The unchanged nonzero
coefficient at $a$ and vanishing below $a$ characterize exactly the same
original least exponent.
\end{proof}

\begin{theorem}[The entire vertex-weight component ideal has a diagram containing the original diagram]
\label{hir:actual-weighted-vertex-ideal-diagram-inclusion}
\uses{hir:actual-weighted-component-retains-initial-exponent}
\lean{Hironaka.HironakaActualWeightedVertexHS.actualWeightedVertexIdeal,Hironaka.HironakaActualWeightedVertexHS.actualInitialDiagram_subset_weightedVertexIdeal}
\leanok
Take any genuine monic vertex family of the unchanged whole original
formal ideal $I$. Span the actual components of its genuine members at
the weights of their actual vertex exponents. The entire resulting
formal ideal has initial diagram containing the entire initial diagram
of $I$, for the same original monomial order.
\end{theorem}
\begin{proof}
Every genuine monic family member has coefficient one at its actual
vertex and vanishes at every smaller exponent. Its component at that
vertex's weight therefore has exactly the same actual initial exponent,
by the preceding coefficient argument. It is a defining generator of
the displayed entire ideal, so the latter's actual initial diagram
contains each original vertex. Every original diagram exponent lies
above an actual original vertex. The genuine diagram of the component
ideal is closed under increasing the exponent, by actual multiplication
with original monomials. Hence it contains every original diagram exponent.
\end{proof}

\begin{theorem}[Every unchanged formal Hilbert--Samuel value of the vertex-weight ideal is at most the original value]
\label{hir:actual-weighted-vertex-ideal-hs-bound}
\uses{hir:actual-weighted-vertex-ideal-diagram-inclusion}
\lean{Hironaka.HironakaActualWeightedVertexHS.genuineFormalHilbertSamuelValue_weightedVertex_le}
\leanok
For the genuine degree-lexicographic vertex family, arbitrary
nonnegative weights and every order $\ell$, the existing whole formal
Hilbert--Samuel value of the actual vertex-weight component ideal is
at most the existing whole formal value of $I$.
\end{theorem}
\begin{proof}
The existing actual standard-monomial theorem identifies each unchanged
whole quotient dimension with the cardinality of the original bounded
exponents outside that ideal's actual initial diagram. The preceding
whole diagram inclusion gives an injection from the bounded complement
for the component ideal into the original bounded complement: retain
the same exponent, and use inclusion contrapositively for its genuine
nonmembership. These are finite actual exponent sets, so their
cardinalities satisfy the asserted inequality. The original standard-
monomial dimension identities give the statement at every arbitrary
order, without an ideal-presentation or Hilbert--Samuel certificate.
\end{proof}

The retained degree is the weight of the genuine initial exponent; it
need not be the least weight present in the series. To identify these
components with actual exceptional restrictions for a positive-dimensional
centre still requires the genuine centre-power compatibility and
analytic coefficient realization. Those producers, and the general
positive-dimensional Bennett theorem, remain separate and open.
